\section{Deployment Instructions}

The application supports two deployment modes: \textbf{Docker Compose} for single-node setups and \textbf{Docker Swarm} (beta) for multi-node, highly available deployments. In both cases, all services run as independent containers and are managed via a centralised script.

\subsection{Installation}

\noindent\textbf{1. Clone the repository.}\par
\begin{lstlisting}[language=bash]
git clone https://github.com/EvenToNight/EvenToNight.git
cd EvenToNight
\end{lstlisting}

\noindent\textbf{2. Configure environment variables.}\par
\begin{lstlisting}[language=bash]
cp .env.template .env
# Edit .env and fill in all required fields
# Note: if using the --no-deps flag, Stripe keys can contain arbitrary values. All fields must still be filled in.
\end{lstlisting}

\subsection{Docker Compose}

\noindent\textbf{Option A: Use pre-built images from ghcr.io (Recommended).}\par
\begin{lstlisting}[language=bash]
# Pull images
./scripts/composeApplication.sh pull

# Pull images with database seeding
./scripts/composeApplication.sh --init-db pull

# Deploy the application
./scripts/composeApplication.sh up -d --wait

# Deploy with database seeding
./scripts/composeApplication.sh --init-db up -d --wait

# Deploy in development mode (with host-mapped ports and dashboards for databases, RabbitMQ and Traefik)
./scripts/composeApplication.sh --init-db --dev up -d --wait
\end{lstlisting}

\noindent\textbf{Option B: Local build.}
Add the \texttt{-{}-build} flag to build services locally instead of using pre-built images. The \texttt{-{}-dev} flag is required as it includes the build instructions:
\begin{lstlisting}[language=bash]
# Build and deploy
./scripts/composeApplication.sh --dev up --build -d --wait

# Build and deploy with seeding
./scripts/composeApplication.sh --init-db --dev up --build -d --wait
\end{lstlisting}

\noindent\textbf{Additional flags.}\par\noindent
\texttt{-{}-no-deps}: Excludes external dependencies (Stripe). The flag can be added to any deploy command to exclude external services:
\begin{lstlisting}[language=bash]
# Deploy without external dependencies
./scripts/composeApplication.sh --no-deps up -d --wait

# Deploy with seeding but without external dependencies
./scripts/composeApplication.sh --init-db --no-deps up -d --wait
\end{lstlisting}

\noindent \textbf{Note:} When using \texttt{-{}-no-deps}, the Stripe keys (\texttt{STRIPE\_SECRET\_KEY}, \texttt{STRIPE\_PUBLISHABLE\_KEY}, \texttt{STRIPE\_WEBHOOK\_SECRET}) in \texttt{.env} can contain arbitrary values.

\noindent\texttt{-{}-project-name}: Overrides the Docker Compose project name (defaults to \texttt{eventonight}):
\begin{lstlisting}[language=bash]
./scripts/composeApplication.sh --project-name myproject up -d --wait
\end{lstlisting}

\noindent\textbf{Stripe configuration.}\label{par:stripe-config} For Stripe payments in a local environment (required only if NOT using \texttt{-{}-no-deps}):
\begin{lstlisting}[language=bash]
./services/payments/scripts/local-webhooks.sh
\end{lstlisting}

\noindent This script must be run to forward Stripe webhooks to the local environment. For more information on using sandbox mode, refer to the \href{https://docs.stripe.com/testing}{Stripe documentation}.

\noindent\textbf{Alternative setup.} Use Gradle to set up the entire environment with seeding and the Stripe listener:
\begin{lstlisting}[language=bash]
./gradlew setupApplicationEnvironment
\end{lstlisting}

\noindent\textbf{Teardown.}\par
\begin{lstlisting}[language=bash]
./scripts/composeApplication.sh down
./scripts/composeApplication.sh down -v   # also remove volumes
./gradlew teardownApplicationEnvironment  # via Gradle
\end{lstlisting}

\subsection{Docker Swarm (Beta)}

\noindent\textbf{Prerequisites.}\par
Initialise the swarm on the manager node:
\begin{lstlisting}[language=bash]
docker swarm init
\end{lstlisting}

\noindent Join additional worker nodes using the token provided by the manager:
\begin{lstlisting}[language=bash]
docker swarm join --token <token> <manager-ip>:2377
\end{lstlisting}

\noindent\textbf{Option A: Use pre-built images from ghcr.io (Recommended).}\par
\begin{lstlisting}[language=bash]
./scripts/deploySwarm.sh
\end{lstlisting}

\noindent The \texttt{-{}-auto-labels} flag can be added to automatically assign placement labels to nodes in a balanced way, without having to configure them manually:
\begin{lstlisting}[language=bash]
./scripts/deploySwarm.sh --auto-labels
\end{lstlisting}

\noindent\textbf{Note:} On the first deploy, database seeding is performed automatically.

\newpage
\noindent\textbf{Option B: Local build.}\par
Build images locally and push them to Docker Hub (multi-arch) so all workers have access to them, then deploy:
\begin{lstlisting}[language=bash]
./scripts/deploySwarm.sh --local --build
\end{lstlisting}

\noindent To override the hostname baked into the frontend image (defaults to \texttt{HOST} from \texttt{.env}, e.g. to use \texttt{localhost} for local testing):
\begin{lstlisting}[language=bash]
./scripts/deploySwarm.sh --local --build --host <hostname> --auto-labels
\end{lstlisting}

\noindent To deploy using already-pushed local images (without rebuilding):
\begin{lstlisting}[language=bash]
./scripts/deploySwarm.sh --local
\end{lstlisting}

\noindent To build and push without deploying:
\begin{lstlisting}[language=bash]
./scripts/deploySwarm.sh --build
\end{lstlisting}

\noindent\textbf{Additional flags.}\par\noindent
\texttt{-{}-no-deps}: Same behaviour as in Docker Compose: excludes external dependencies. Stripe configuration applies equally (see p.~\pageref{par:stripe-config}).
\begin{lstlisting}[language=bash]
./scripts/deploySwarm.sh --no-deps
\end{lstlisting}

\noindent\texttt{-{}-stack-name}: Overrides the stack name (defaults to \texttt{eventonight-swarm}):
\begin{lstlisting}[language=bash]
./scripts/deploySwarm.sh --stack-name mystack
\end{lstlisting}

\noindent\textbf{Recovery.}\par
If some services are not fully running after a deploy, the recovery script force-updates only the failing ones:
\begin{lstlisting}[language=bash]
./scripts/swarmRecover.sh [STACK_NAME]
\end{lstlisting}

\noindent To check the current status of all services without recovering:
\begin{lstlisting}[language=bash]
./scripts/swarmRecover.sh [STACK_NAME] --status
\end{lstlisting}

\noindent \texttt{STACK\_NAME} defaults to \texttt{eventonight-swarm}.

\noindent\textbf{Teardown.}\par
\begin{lstlisting}[language=bash]
./scripts/deploySwarm.sh --stop
./scripts/deploySwarm.sh --stop --remove-volumes
./scripts/deploySwarm.sh --remove-local-images  # remove test images from Docker Hub
\end{lstlisting}

Alternatively, the application is already running in production at \href{https://eventonight.site/it}{EvenToNight}.
